\oldpage[III]{157}
On peut considérer que les isomorphismes (6.6.7, $t$)) définissent des systèmes
inductifs de foncteurs spectraux sur l'ensemble filtrant des couples
$(\mathfrak U^{(1)},\mathfrak U^{(2)})$ de recouvrements ouverts affines finis
de $X^{(1)}$ et $X^{(2)}$. Nous désignerons par
\[
  {}^{(t)}E
  (X^{(1)},X^{(2)};
   \mathscr P_\bullet^{(1)},\mathscr P_\bullet^{(2)})
  \quad\text{ou}\quad
  {}^{(t)}E^S
  (X^{(1)},X^{(2)};
   \mathscr P_\bullet^{(1)},\mathscr P_\bullet^{(2)})
\]
la limite inductive de ce système, par
$\mathbf{Tor}_\bullet^S
(X^{(1)},X^{(2)};\mathscr P_\bullet^{(1)},\mathscr P_\bullet^{(2)})$
l'aboutissement de ce foncteur spectral, que nous appellerons l'\emph{hypertor
global} des deux complexes $\mathscr P_\bullet^{(1)}$ et
$\mathscr P_\bullet^{(2)}$; si $\mathscr P_\bullet^{(1)}$ et
$\mathscr P_\bullet^{(2)}$ sont réduits à leurs termes de degré 0,
$\mathscr F^{(1)}$ et $\mathscr F^{(2)}$, nous écrirons
\[
  \operatorname{Tor}_\bullet^S
  (X^{(1)},X^{(2)};\mathscr F^{(1)},\mathscr F^{(2)}),
\]
et conformément aux conventions générales,
$
  \operatorname{Tor}_q^S
  (X^{(1)},X^{(2)};\mathscr P_\bullet^{(1)},\mathscr P_\bullet^{(2)})
$
sera donc le bicomplexe des
$
  \operatorname{Tor}_q^S
  (X^{(1)},X^{(2)};\mathscr P_h^{(1)},\mathscr P_k^{(2)})$.
\end{env}

\begin{env}[6.6.8]
Les hypothèses étant celles de (6.6.1), considérons maintenant deux
$S$-morphismes $f_i:X^{(i)}\to Y^{(i)}$, où
$Y^{(i)}=\operatorname{Spec}(B_i)$ est un $S$-\emph{schéma affine}, $B_i$
étant donc une $A$-algèbre ($i=1,2$); cela définit donc un $A$-homomorphisme
$B_i\to\Gamma(X^{(i)},\mathscr O_{X^{(i)}})$ (\textbf{I}, 2.2.4), et par suite chacun
des $L_{\bullet\bullet}^{(i)}$ définis dans (6.6.2) est un bicomplexe de
$B_i$-modules; on en conclut que
$L_{\bullet\bullet}^{(1)}\otimes_A L_{\bullet\bullet}^{(2)}$ est un
quadricomplexe de $(B_1\otimes_A B_2)$-modules, et ses six foncteurs spectraux
d'hyperhomologie peuvent donc être considérés comme prenant leurs valeurs
dans la catégorie des suites spectrales de $(B_1\otimes_A B_2)$-\emph{modules}. Si
l'on pose
$Y=Y^{(1)}\times_S Y^{(2)}=\operatorname{Spec}(B_1\otimes_A B_2)$, on peut
considérer les $\mathscr O_Y$-Modules quasi-cohérents associés à ces modules
(\textbf{I}, 1.3.4); nous noterons
${}^{(t)}\mathscr E
(f_1,f_2;\mathscr P_\bullet^{(1)},\mathscr P_\bullet^{(2)})$
(pour $t=a$, $b$, $a'$, $b'$, $c$ ou $d$) les six suites spectrales
$
  \bigl({}^{(t)}E
  (X^{(1)},X^{(2)};
   \mathscr P_\bullet^{(1)},\mathscr P_\bullet^{(2)})\bigr)^{\sim}$
de $\mathscr O_Y$-Modules, et
$\boldsymbol{\mathscr{T}\!or}_\bullet^S
(f_1,f_2;\mathscr P_\bullet^{(1)},\mathscr P_\bullet^{(2)})$ leur
aboutissement commun
$
  \bigl(\mathbf{Tor}_\bullet^S
  (X_\bullet^{(1)},X^{(2)};
   \mathscr P_\bullet^{(1)},\mathscr P_\bullet^{(2)})\bigr)^{\sim}$.
On le notera
$\mathscr{T}\!or_\bullet^S(f_1,f_2;\mathscr F^{(1)},\mathscr F^{(2)})$
lorsque $\mathscr P_\bullet^{(i)}$ est réduit à son terme de degré 0,
$\mathscr F^{(i)}$ ($i=1,2$).
\end{env}

\subsection{Foncteurs hypertor globaux de complexes de Modules quasi-cohérents
et suites spectrales de Künneth : cas général.}

\begin{env}[6.7.1]
Nous allons maintenant généraliser les définitions de (6.6.8) au cas où $S$
est un préschéma quelconque, $X^{(i)}$, $Y^{(i)}$ des $S$-préschémas et
$f_i:X^{(i)}\to Y^{(i)}$ des morphismes \emph{séparés} et \emph{quasi-compacts}. Il
s'agit alors, pour tout couple de complexes $\mathscr P_\bullet^{(i)}$ de
$\mathscr O_{X^{(i)}}$-Modules quasi-cohérents, limités inférieurement,
($i=1,2$), de définir pour tout $n$ un $\mathscr O_Y$-Module quasi-cohérent
$\boldsymbol{\mathscr{T}\!or}_n^S
(f_1,f_2;\mathscr P_\bullet^{(1)},\mathscr P_\bullet^{(2)})$ ainsi que 6
foncteurs spectraux, se réduisant aux définitions de (6.6.8) lorsque $S$,
$Y^{(1)}$ et $Y^{(2)}$ sont affines (on pose
$Y=Y^{(1)}\times_S Y^{(2)}$).
Supposons d'abord $S=\operatorname{Spec}(A)$ \emph{affine}, mais $Y^{(1)}$
et $Y^{(2)}$ quelconques; soit $W^{(i)}$ un ouvert affine de $Y^{(i)}$;
$f_i^{-1}(W^{(i)})$ est alors un $S$-schéma quasi-compact,
$W=W^{(1)}\times_S W^{(2)}$ un ouvert affine de $Y$; soit
$f'_i:f_i^{-1}(W^{(i)})\to W^{(i)}$ la restriction de $f_i$, et
$\mathscr P'_\bullet{}^{(i)}$ la restriction
$\mathscr P_\bullet^{(i)}|f_i^{-1}(W^{(i)})$ ($i=1,2$). On a alors d'après
(6.6.8) les suites spectrales
${}^{(t)}\mathscr E
(f'_1,f'_2;\mathscr P'_\bullet{}^{(1)},\mathscr P'_\bullet{}^{(2)})$
de $(\mathscr O_Y|W)$-Modules quasi-cohérents, et il s'agit de vérifier
qu'elles satisfont aux conditions de recollement ($\mathbf{0}_{\mathrm I}$, 3.3.1).
On est aussitôt ramené au cas où $Y^{(i)}=\operatorname{Spec}(B_i)$ est affine
et où $W^{(i)}=D(g_i)$, où $g_i\in B_i$, de sorte que
$W=D(g_1\otimes g_2)$
\oldpage[III]{158}
dans $Y=\operatorname{Spec}(B_1\otimes_A B_2)$ (\textbf{II}, 4.3.2.1); si
$X'^{(i)}=f_i^{-1}(W^{(i)})$, il s'agit d'établir un isomorphisme canonique
de foncteurs spectraux
\[
  {}^{(t)}E
  (X'^{(1)},X'^{(2)};
   \mathscr P'_\bullet{}^{(1)},\mathscr P'_\bullet{}^{(2)})
  \xrightarrow{\sim}
  {}^{(t)}E
  (X^{(1)},X^{(2)};
   \mathscr P_\bullet^{(1)},\mathscr P_\bullet^{(2)})
  \otimes_B B_g
  \tag{6.7.1.1}\label{III.6.7.1.1.fr}
\]
où l'on a posé $B=B_1\otimes_A B_2$ et $g=g_1\otimes g_2$. Pour cela,
partons de recouvrements ouverts affines finis $\mathfrak U^{(i)}$ de
$X^{(i)}$ ($i=1,2$), et soit $\mathfrak U'^{(i)}$ la trace de
$\mathfrak U^{(i)}$ sur $X'^{(i)}$, qui est encore formé d'ouverts affines
(\textbf{I}, 5.5.10); de façon précise, on a
\[
  C^\bullet(\mathfrak U'^{(i)},\mathscr P'_\bullet{}^{(i)})
  =C^\bullet(\mathfrak U^{(i)},\mathscr P_\bullet^{(i)})
   \otimes_{B_i}(B_i)_{g_i}.
\]
Si on pose
$L'_{\bullet\bullet}{}^{(i)}
=C^\bullet(\mathfrak U'^{(i)},\mathscr P'_\bullet{}^{(i)})$, on a donc
$
  L'_{\bullet\bullet}{}^{(1)}\otimes_A L'_{\bullet\bullet}{}^{(2)}
  =\bigl(L_{\bullet\bullet}^{(1)}\otimes_{B_1}(B_1)_{g_1}\bigr)
   \otimes_A
   \bigl(L_{\bullet\bullet}^{(2)}\otimes_{B_2}(B_2)_{g_2}\bigr)$;
comme on a $B_g=(B_1)_{g_1}\otimes_A(B_2)_{g_2}$ à un isomorphisme
canonique près, on a, à un isomorphisme canonique près,
$
  L'_{\bullet\bullet}{}^{(1)}\otimes_A L'_{\bullet\bullet}{}^{(2)}
  =(L_{\bullet\bullet}^{(1)}\otimes_A L_{\bullet\bullet}^{(2)})
    \otimes_B B_g$.
Si $M_{\bullet\bullet\bullet}^{(i)}$ est une résolution projective de
Cartan-Eilenberg de $L_{\bullet\bullet}^{(i)}$, qu'on peut supposer formée de
$B_i$-modules, il résulte du fait que $(B_i)_{g_i}$ est \emph{plat} sur $B_i$
que
$M'_{\bullet\bullet\bullet}{}^{(i)}
=M_{\bullet\bullet\bullet}^{(i)}\otimes_{B_i}(B_i)_{g_i}$ est une résolution
projective de Cartan-Eilenberg du bicomplexe
$L'_{\bullet\bullet}{}^{(i)}$; en outre, on a
\[
  M'_{\bullet\bullet\bullet}{}^{(1)}\otimes_A M'_{\bullet\bullet\bullet}{}^{(2)}
  =(M_{\bullet\bullet\bullet}^{(1)}\otimes_A M_{\bullet\bullet\bullet}^{(2)})
    \otimes_B B_g.
\]
L'isomorphisme cherché (6.7.1.1) résulte alors aussitôt des définitions de
l'hyperhomologie d'un bicomplexe, et de l'exactitude du foncteur
$G\otimes_B B_g$ en le $B$-module $G$.
\end{env}

\begin{env}[6.7.2]
Supposons maintenant $S$ quelconque, et soient
$u_i:Y^{(i)}\to S$ les morphismes structuraux ($i=1,2$). Soit $(S_\alpha)$ un
recouvrement ouvert affine de $S$, posons
$Y_\alpha^{(i)}=u_i^{-1}(S_\alpha)$,
$X_\alpha^{(i)}=f_i^{-1}(Y_\alpha^{(i)})$, et soit
$f_{i\alpha}:X_\alpha^{(i)}\to Y_\alpha^{(i)}$ la restriction de $f_i$, qui
est un morphisme séparé et quasi-compact. Les
$Y_\alpha=Y_\alpha^{(1)}\times_{S_\alpha}Y_\alpha^{(2)}$ forment un
recouvrement ouvert de $Y$, et sur chaque $Y_\alpha$ sont définis par (6.7.1)
des foncteurs spectraux
\[
  {}^{(t)}\mathscr E^{S_\alpha}
  \bigl(f_{1\alpha},f_{2\alpha};
    \mathscr P_\bullet^{(1)}|X_\alpha^{(1)},
    \mathscr P_\bullet^{(2)}|X_\alpha^{(2)}\bigr);
\]
il s'agit encore de montrer que ces foncteurs vérifient les conditions de
recollement.
On se ramène aussitôt à la situation suivante :
$S=\operatorname{Spec}(A)$ est affine, $S'=D(h)$, où $h\in A$, et
$u_i(Y^{(i)})\subset S'$; on peut en outre supposer
$Y^{(i)}=\operatorname{Spec}(B_i)$ affine; il s'agit de définir des
isomorphismes canoniques
\[
  {}^{(t)}E^S
  (X^{(1)},X^{(2)};
   \mathscr P_\bullet^{(1)},\mathscr P_\bullet^{(2)})
  \xrightarrow{\sim}
  {}^{(t)}E^{S'}
  (X^{(1)},X^{(2)};
   \mathscr P_\bullet^{(1)},\mathscr P_\bullet^{(2)}).
  \tag{6.7.2.1}\label{III.6.7.2.1.fr}
\]

Or, avec les notations de (6.6.2), les $L_{\bullet\bullet}^{(i)}$ sont formés
de $A_h$-modules, et l'on a donc
$L_{\bullet\bullet}^{(1)}\otimes_{A_h}L_{\bullet\bullet}^{(2)}
=L_{\bullet\bullet}^{(1)}\otimes_A L_{\bullet\bullet}^{(2)}$ à un
isomorphisme canonique près; comme on peut prendre une résolution projective
de Cartan-Eilenberg $M_{\bullet\bullet\bullet}^{(i)}$ de
$L_{\bullet\bullet}^{(i)}$ formée de $A_h$-modules, cela donne aussitôt
l'isomorphisme canonique cherché.
\end{env}

Nous avons en résumé démontré le

\begin{theorem}[6.7.3]
Soient $S$ un préschéma, $f_i:X_i\to Y_i$ un $S$-morphisme séparé et
quasi-compact de $S$-préschémas, $\mathscr P_\bullet^{(i)}$ un complexe de
$\mathscr O_{X^{(i)}}$-Modules quasi-cohérents limité inférieurement
($i=1,2$); on pose $Y=Y^{(1)}\times_S Y^{(2)}$. Il existe un
bi-$\partial$-foncteur
  $\boldsymbol{\mathscr{T}\!or}_\bullet^S
(f_1,f_2;\mathscr P_\bullet^{(1)},\mathscr P_\bullet^{(2)})$
\oldpage[III]{159}
à valeurs dans la catégorie des $\mathscr O_Y$-Modules quasi-cohérents, tel
que si $V^{(i)}$ est un ouvert affine de $Y^{(i)}$ ($i=1,2$) et
$V=V^{(1)}\times_S V^{(2)}$, on ait
\[
  \boldsymbol{\mathscr{T}\!or}_\bullet^S
  (f_1,f_2;\mathscr P_\bullet^{(1)},\mathscr P_\bullet^{(2)})|V
  =\Bigl(\mathbf{Tor}_\bullet^S
    \bigl(f_1^{-1}(V^{(1)}),f_2^{-1}(V^{(2)});
      \mathscr P_\bullet^{(1)}|f_1^{-1}(V^{(1)}),
      \mathscr P_\bullet^{(2)}|f_2^{-1}(V^{(2)})\bigr)\Bigr)^{\sim}.
\]

Ce bifoncteur est l'aboutissement de six foncteurs spectraux biréguliers
\[
  {}^{(t)}\mathscr E
  (f_1,f_2;\mathscr P_\bullet^{(1)},\mathscr P_\bullet^{(2)})
  \qquad(t=a,b,a',b',c,d)
\]
dont les termes $E_2$ sont donnés par
\[
\begin{aligned}
  {}^{(a)}\mathscr E_{pq}^2
  &=\bigoplus_{q_1+q_2=q}
    \mathscr{T}\!or_p^S
    \bigl(\mathcal H^{-q_1}(f_1,\mathscr P_\bullet^{(1)}),
          \mathcal H^{-q_2}(f_2,\mathscr P_\bullet^{(2)})\bigr),\\
  {}^{(b)}\mathscr E_{pq}^2
  &=\mathscr H^{-p}\bigl(f_1\times_S f_2,
    \boldsymbol{\mathscr{T}\!or}_q^S
    (\mathscr P_\bullet^{(1)},\mathscr P_\bullet^{(2)})\bigr),\\
  {}^{(a')}\mathscr E_{pq}^2
  &=\bigoplus_{q_1+q_2=q}
    \boldsymbol{\mathscr{T}\!or}_p^S
    \bigl(\mathscr H^{-q_1}(f_1,\mathscr P_\bullet^{(1)}),
          \mathscr H^{-q_2}(f_2,\mathscr P_\bullet^{(2)})\bigr),\\
  {}^{(b')}\mathscr E_{pq}^2
  &=\mathcal H^{-p}\bigl(f_1\times_S f_2,
    \mathscr{T}\!or_q^S
    (\mathscr P_\bullet^{(1)},\mathscr P_\bullet^{(2)})\bigr),\\
  {}^{(c)}\mathscr E_{pq}^2
  &=\bigoplus_{q_1+q_2=q}
    \mathscr{T}\!or_p^S
    \bigl(f_1,f_2;
      \mathscr H_{q_1}(\mathscr P_\bullet^{(1)}),
      \mathscr H_{q_2}(\mathscr P_\bullet^{(2)})\bigr),\\
  {}^{(d)}\mathscr E_{pq}^2
  &=\mathscr H_p\bigl(\mathscr{T}\!or_q^S
    (f_1,f_2;\mathscr P_\bullet^{(1)},\mathscr P_\bullet^{(2)})\bigr).
\end{aligned}
\]
\end{theorem}
\noindent On dit que les suites spectrales $(a)$ et $(b)$ sont les \emph{suites
spectrales de Künneth}.

On notera que les suites spectrales $(a)$ et $(a')$ (resp. $(b)$ et $(b')$)
sont identiques lorsque $\mathscr P_\bullet^{(1)}$ et
$\mathscr P_\bullet^{(2)}$ se réduisent à leurs termes de degré 0; dans ce cas,
les suites $(c)$ et $(d)$ sont dégénérées et sont donc sans intérêt.

\begin{remark}[6.7.4]
Les hypertor globaux que nous avons définis ci-dessus comprennent comme cas
particuliers, à la fois les \emph{Modules d'hypercohomologie} définis dans
(6.2.1) et les \emph{hypertor locaux} définis dans (6.5.3). Montrons que l'on
a, pour tout morphisme $f:X\to Y$ \emph{quasi-compact} et \emph{séparé} et tout
complexe $\mathscr P_\bullet$ de $\mathscr O_X$-Modules quasi-cohérents,
limité inférieurement, un isomorphisme canonique de $\partial$-foncteurs en
$\mathscr P_\bullet$
\[
  \boldsymbol{\mathscr{T}\!or}_n^Y(f,1_Y;\mathscr P_\bullet,\mathscr O_Y)
  \xrightarrow{\sim}
  \mathcal H^{-n}(f,\mathscr P_\bullet)
  \qquad\text{(pour tout $n\in\mathbf Z$)}.
  \tag{6.7.4.1}\label{III.6.7.4.1.fr}
\]

En effet, les méthodes de recollement de (6.7.2) ramènent aussitôt au cas où
$Y$ est affine; on peut alors, en vertu de (6.2.2), calculer les deux membres
de (6.7.4.1) à l'aide d'un même recouvrement fini $\mathfrak U$ de $Y$ par des
ouverts affines, et (pour le premier membre) du recouvrement de $Y$ formé de
$Y$ lui-même; avec les notations de (6.6.2), le bicomplexe
$L_{\bullet\bullet}^{(2)}$ est alors réduit à son terme de degrés $(0,0)$,
égal à $A$, et la conclusion résulte de (\textbf{0}, 11.7.5). Pour une
généralisation de ce résultat, voir (6.7.7); mais on notera que lorsque dans
le premier membre de (6.7.4.1), on remplace $\mathscr O_Y$ par un
$\mathscr O_Y$-Module quasi-cohérent \emph{quelconque} $\mathscr F$, on n'a
plus en général un isomorphisme avec
$\mathcal H^{-n}(f,\mathscr P_\bullet\otimes_{\mathscr O_Y}\mathscr F)$, bien
que, dans le calcul précédent, le bicomplexe
$L_{\bullet\bullet}^{(1)}\otimes_A L_{\bullet\bullet}^{(2)}$ s'identifie
encore au bicomplexe
$C^\bullet(\mathfrak U,\mathscr P_\bullet\otimes_Y\mathscr F)$.

D'autre part, on a un isomorphisme canonique de bi-$\partial$-foncteurs
\[
  \boldsymbol{\mathscr{T}\!or}_\bullet^S
  (1_{X^{(1)}},1_{X^{(2)}};
   \mathscr P_\bullet^{(1)},\mathscr P_\bullet^{(2)})
  \xrightarrow{\sim}
  \boldsymbol{\mathscr{T}\!or}_\bullet^S
  (\mathscr P_\bullet^{(1)},\mathscr P_\bullet^{(2)}).
  \tag{6.7.4.2}\label{III.6.7.4.2.fr}
\]

En effet, on se ramène encore, par (6.7.1) et (6.7.2), au cas où $S$ et les
$X^{(i)}$ sont affines; dans le calcul du premier membre de (6.7.4.2), on peut
alors prendre pour
\oldpage[III]{160}
recouvrement $\mathfrak U^{(i)}$ la famille réduite au seul élément $X^{(i)}$,
de sorte qu'avec les notations de (6.6.2), $L_{\bullet\bullet}^{(i)}$ se réduit
à $\Gamma(X^{(i)},\mathscr P_\bullet^{(i)})$ (considéré comme bicomplexe dont
les termes de premier degré $\ne0$ sont nuls), et l'égalité des deux membres
de (6.7.4.2) résulte de (6.4.1.1) et (6.3.1).
\end{remark}

\begin{proposition}[6.7.5]
Soit $u:\mathscr P_\bullet^{(1)}\to\mathscr Q_\bullet^{(1)}$ un
homomorphisme de complexes de $\mathscr O_{X^{(1)}}$-Modules quasi-cohérents,
limités inférieurement, tel que l'homomorphisme
\[
  \mathscr H_\bullet(u):
  \mathscr H_\bullet(\mathscr P_\bullet^{(1)})
  \longrightarrow
  \mathscr H_\bullet(\mathscr Q_\bullet^{(1)})
\]
déduit de $u$ soit un isomorphisme. Alors les homomorphismes
\[
  {}^{(t)}\mathscr E
  (f_1,f_2;\mathscr P_\bullet^{(1)},\mathscr P_\bullet^{(2)})
  \longrightarrow
  {}^{(t)}\mathscr E
  (f_1,f_2;\mathscr Q_\bullet^{(1)},\mathscr P_\bullet^{(2)})
\]
déduits de $u$ sont des isomorphismes pour $t=a$, $t=b$ et $t=c$.
\end{proposition}

\begin{proof}
L'assertion relative à la suite spectrale $(c)$ résulte de ce que cette suite
est birégulière et de ce que l'homomorphisme considéré est un isomorphisme pour
les termes $E_2$ par hypothèse (\textbf{0}, 11.1.5). Ceci montre déjà que
$
  \boldsymbol{\mathscr{T}\!or}_\bullet^S
  (f_1,f_2;\mathscr P_\bullet^{(1)},\mathscr P_\bullet^{(2)})
  \longrightarrow
  \boldsymbol{\mathscr{T}\!or}_\bullet^S
  (f_1,f_2;\mathscr Q_\bullet^{(1)},\mathscr P_\bullet^{(2)})
$
est un isomorphisme. Appliquant les relations (6.7.4.1) et (6.7.4.2) on voit
d'abord que les homomorphismes
$
  \mathcal H^{-n}(f_1,\mathscr P_\bullet^{(1)})
  \longrightarrow
  \mathcal H^{-n}(f_1,\mathscr Q_\bullet^{(1)})
  $ et $
  \boldsymbol{\mathscr{T}\!or}_n^S
  (\mathscr P_\bullet^{(1)},\mathscr P_\bullet^{(2)})
  \longrightarrow
  \boldsymbol{\mathscr{T}\!or}_n^S
  (\mathscr Q_\bullet^{(1)},\mathscr P_\bullet^{(2)})
$
déduits de $u$ sont des isomorphismes. L'assertion relative aux suites $(a)$
et $(b)$ résulte alors de ce que ces suites sont birégulières (6.7.3) et que
les homomorphismes considérés sont bijectifs pour les termes $E_2$
(\textbf{0}, 11.1.5).

Notons d'autre part que, si
$u:\mathscr P_\bullet^{(1)}\to\mathscr Q_\bullet^{(1)}$ est un
\emph{homotopisme}, on en déduit des isomorphismes canoniques
$
  {}^{(t)}\mathscr E
  (f_1,f_2;\mathscr P_\bullet^{(1)},\mathscr P_\bullet^{(2)})
  \longrightarrow
  {}^{(t)}\mathscr E
  (f_1,f_2;\mathscr Q_\bullet^{(1)},\mathscr P_\bullet^{(2)})
$
pour les six suites spectrales. En effet, si $S$ et les $Y^{(i)}$ sont affines,
on déduit de $u$ un homotopisme de bicomplexes
$
  C^\bullet(\mathfrak U^{(1)},\mathscr P_\bullet^{(1)})
  \longrightarrow
  C^\bullet(\mathfrak U^{(1)},\mathscr Q_\bullet^{(1)})$,
et la proposition résulte de la théorie générale de l'hyperhomologie
(\textbf{0}, 11.3.2); le passage au cas général se fait par recollement,
en utilisant le fait que, d'un homotopisme de complexes, on déduit un
homotopisme de résolutions projectives de Cartan-Eilenberg de ces complexes
(M, XVII, 1.2).
\end{proof}

\begin{proposition}[6.7.6]
Supposons que le complexe $\mathscr P_\bullet^{(1)}$ ou le complexe
$\mathscr P_\bullet^{(2)}$ soit formé de Modules $S$-plats (les deux complexes
étant limités inférieurement). On a alors un isomorphisme canonique de
bi-$\partial$-foncteurs
\[
  \boldsymbol{\mathscr{T}\!or}_n^S
  (f_1,f_2;\mathscr P_\bullet^{(1)},\mathscr P_\bullet^{(2)})
  \xrightarrow{\sim}
  \mathcal H^{-n}\bigl(f_1\times_S f_2,
    \mathscr P_\bullet^{(1)}\otimes_S\mathscr P_\bullet^{(2)}\bigr).
  \tag{6.7.6.1}\label{III.6.7.6.1.fr}
\]
\end{proposition}

\begin{proof}
Supposons d'abord $S$, $Y^{(1)}$ et $Y^{(2)}$ affines, de sorte qu'on est
dans la situation de (6.6.2), dont nous conservons les notations. Supposons par
exemple que $\mathscr P_\bullet^{(1)}$ soit formé de Modules $S$-plats, et
calculons l'hypertor en utilisant la remarque (6.6.6) : c'est donc l'homologie
de $L_{\bullet\bullet}^{(1)}\otimes_A M_{\bullet\bullet\bullet}^{(2)}$, où
$M_{\bullet\bullet\bullet}^{(2)}$ est une résolution projective de
Cartan-Eilenberg de $L_{\bullet\bullet}^{(2)}$, au sens de
(\textbf{0}, 11.7.1). D'autre part, les modules
$L_{ij}^{(1)}$ sont \emph{plats} sur $A$ en vertu de l'hypothèse (1.4.15.1);
on déduit alors de (\textbf{0}, 11.7.5) un isomorphisme canonique
\[
  \mathbf{Tor}_\bullet^S
  (\mathfrak U^{(1)},\mathfrak U^{(2)};
   \mathscr P_\bullet^{(1)},\mathscr P_\bullet^{(2)})
  \xrightarrow{\sim}
  \mathscr H_\bullet
  (L_{\bullet\bullet}^{(1)}\otimes_A L_{\bullet\bullet}^{(2)}).
  \tag{6.7.6.2}\label{III.6.7.6.2.fr}
\]

On a d'autre part un homomorphisme naturel de bicomplexes de
$L_{\bullet\bullet}^{(1)}\otimes_A L_{\bullet\bullet}^{(2)}$ dans
$C^\bullet(\mathfrak U,\mathscr Q_\bullet)$, où $\mathfrak U$ est le
recouvrement de $Z=X^{(1)}\times_S X^{(2)}$ par les ouverts affines
\oldpage[III]{161}
$U_\alpha^{(1)}\times_S U_\beta^{(2)}$ et
$\mathscr Q_\bullet
=\mathscr P_\bullet^{(1)}\otimes_S\mathscr P_\bullet^{(2)}$ (considéré comme
complexe simple pour le degré total); en effet, la définition de cet
homomorphisme a en substance été donnée au cours du calcul de la suite $(b')$
dans (6.6.4), pour $q=0$; il suffit simplement (en gardant les notations de
(6.6.4)) de tenir compte de ce qu'il y a d'une part un homomorphisme naturel
du complexe $N^{\bullet k}$ dans le complexe
$C^\bullet(\Phi^{(1)},\Phi^{(2)};\mathscr S_k)$
(\textbf{0}, 11.8.5), d'autre part un homomorphisme naturel de ce dernier
complexe dans le complexe
$P^\bullet(\Phi^{(1)},\Phi^{(2)};\mathscr S_k)$
(\textbf{0}, 11.8.6), et enfin un homomorphisme naturel de ce dernier
complexe de cochaînes dans le sous-complexe des cochaînes alternées
(\textbf{0}, 11.8.7). Par ailleurs, l'homomorphisme de bicomplexes ainsi
défini donne un isomorphisme en homologie, comme on l'a vu en (6.6.4); on a
donc, en composant avec (6.7.6.2), obtenu un isomorphisme
\[
  \mathbf{Tor}_n^S
  (\mathfrak U^{(1)},\mathfrak U^{(2)};
   \mathscr P_\bullet^{(1)},\mathscr P_\bullet^{(2)})
  \xrightarrow{\sim}
  \mathbf H^{-n}
  (\mathfrak U,\mathscr P_\bullet^{(1)}\otimes_S\mathscr P_\bullet^{(2)}).
  \tag{6.7.6.3}\label{III.6.7.6.3.fr}
\]

Il faut ensuite prouver que l'isomorphisme ainsi défini ne dépend pas des
recouvrements ouverts choisis (le second membre de (6.7.6.3) étant
canoniquement isomorphe à
$\mathbf H^{-n}
(X^{(1)}\times_S X^{(2)},
 \mathscr P_\bullet^{(1)}\otimes_S\mathscr P_\bullet^{(2)})$ par (6.2.2));
cela se fait à l'aide de (6.6.7) en remarquant (avec les notations de (6.6.7))
que l'on a un diagramme commutatif à homotopismes près
\[
\xymatrix@C=18pt@R=28pt{
  C^\bullet(\mathfrak U^{(1)},\mathscr P_\bullet^{(1)})
  \otimes_A
  C^\bullet(\mathfrak U^{(2)},\mathscr P_\bullet^{(2)})
  \ar[r]\ar[d]
  & C^\bullet(\mathfrak U,\mathscr Q_\bullet)\ar[d]
  \\
  C^\bullet(\mathfrak V^{(1)},\mathscr P_\bullet^{(1)})
  \otimes_A
  C^\bullet(\mathfrak V^{(2)},\mathscr P_\bullet^{(2)})
  \ar[r]
  & C^\bullet(\mathfrak V,\mathscr Q_\bullet)
}
\]
où les flèches horizontales sont les homomorphismes définis ci-dessus. Enfin,
il faut passer au cas général par recollement, ce qui se fait sans difficulté
comme dans (6.7.1) et (6.7.2); nous laissons les détails au lecteur.
\end{proof}

\begin{proposition}[6.7.7]
Supposons que $\mathscr P_\bullet^{(1)}$ et
$\mathscr P_\bullet^{(2)}$ soient limités inférieurement, et que tous les
Modules $\mathcal H^{-n}(f_1,\mathscr P_\bullet^{(1)})$ \textup{ou} tous les Modules
$\mathcal H^{-n}(f_2,\mathscr P_\bullet^{(2)})$ soient $S$-plats. On a alors un
isomorphisme canonique de bi-$\partial$-foncteurs ($n$ parcourant $\mathbf Z$)
\[
  \boldsymbol{\mathscr{T}\!or}_n^S
  (f_1,f_2;\mathscr P_\bullet^{(1)},\mathscr P_\bullet^{(2)})
  \xrightarrow{\sim}
  \bigoplus_{q_1+q_2=n}
  \mathcal H^{-q_1}(f_1,\mathscr P_\bullet^{(1)})
  \otimes_S
  \mathcal H^{-q_2}(f_2,\mathscr P_\bullet^{(2)}).
  \tag{6.7.7.1}\label{III.6.7.7.1.fr}
\]
\end{proposition}

\begin{proof}
En effet, vu (6.5.8), la suite spectrale $(a)$ de (6.7.3) est dégénérée, et
la proposition résulte aussitôt de (\textbf{0}, 11.1.6), cette suite
étant birégulière (6.7.3).
\end{proof}

\begin{theorem}[6.7.8]
Supposons que : $1^\circ$ les complexes $\mathscr P_\bullet^{(1)}$ et
$\mathscr P_\bullet^{(2)}$ soient limités inférieurement; $2^\circ$ le
complexe $\mathscr P_\bullet^{(1)}$ \textup{ou} le complexe
$\mathscr P_\bullet^{(2)}$ soit formé de Modules $S$-plats; $3^\circ$ tous les
